\documentclass{article}
\usepackage[UTF8]{ctex}
\usepackage{amsmath,amssymb,amsfonts}
\usepackage{enumitem}
\usepackage{geometry}
\geometry{a4paper, margin=1in}
% === 错题本排版优化（不参与题目标识） ===
\setlength{\parindent}{0pt}
\setlength{\parskip}{0.35em plus 0.10em minus 0.08em}
\setlength{\emergencystretch}{3em}
\setlist{itemsep=0.25em, topsep=0.40em, parsep=0pt}
\allowdisplaybreaks[2]
% === 错题本排版优化结束 ===
\begin{document}
\section*{北京邮电大学2024年高等代数真题}
\section*{12．北京邮电大学2024年高等代数试卷}
一．证明：在 \(\mathbb{Q}[x]\) 中，如果 \((x+1) \mid f\left(x^{2 k+1}\right)\) ，那么 \(\left(x^{2 k+1}+1\right) \mid f\left(x^{2 k+1}\right)\) ，其中 \(k\) 为任意自然数。

二．已知矩阵

\[
A=\left(\begin{array}{lll}
1 & 1 & 0 \\
0 & 0 & 1 \\
0 & 1 & 0
\end{array}\right)
\]

1．证明当 \(n \geq 3\) 时，有 \(A^{n}=A^{n-2}+A^{2}-E\) ，其中 \(E\) 为3阶单位矩阵。\\

2 ．求 \(A^{100}\) ．

三．当

\[
A^{-1}=\left(\begin{array}{lll}
1 & 1 & 1 \\
1 & 2 & 1 \\
1 & 1 & 3
\end{array}\right)
\]

求 \(A^{*},\left|A^{*}\right|, \sum_{j=1}^{3} \sum_{i=1}^{3} A_{i j}\) ，其中 \(A^{*}\) 是 \(A\) 的伴随矩阵，\(A_{i j}\) 是 \(|A|\) 中第 \(i\) 行第 \(j\) 列的代数余子式。

四．已知 \(\alpha_{1}=(1,2,3), \alpha_{2}=(3,-1,2), \alpha_{3}=(2,3, c)\) ，当 \(c\) 为何值时，\(\alpha_{1}, \alpha_{2}, \alpha_{3}\) 线性相关？当 \(c\) 为何值时，\\
\(\alpha_{1}, \alpha_{2}, \alpha_{3}\) 线性无关？

五。设3阶实对称矩阵 \(A\) 的特征值是 \(1,2,3, \alpha_{1}=(-1,-1,1)^{\prime}, \alpha_{2}=(1,-2,-1)^{\prime}\) 分别为 \(A\) 对应于 1,2的特征向量．

\begin{itemize}

  \item[1．] 求 \(A\) 对应于特征值3的特征向量。

  \item[2．] 求矩阵 \(A\) 。

  \item[3．] 设 \(\beta=(2,1,0)^{\prime}\) ，计算 \(A^{n} \beta\) ．
\end{itemize}

六。已知实二次型 \(f\left(x_{1}, x_{2}, x_{3}\right)=a x_{1}^{2}+2 x_{2}^{2}-2 x_{3}^{2}+2 b x_{1} x_{3}\)（其中 \(b>0\) ），它的矩阵 \(A\) 的特征值之和为1，特征值之积为－12。

\begin{itemize}

  \item[1．] 求 \(a, b\) 的值。

  \item[2．] 求正交变换，把上述二次型 \(f\left(x_{1}, x_{2}, x_{3}\right)\) 化为标准形．
\end{itemize}

七．设实矩阵 \(A=\left(a_{i j}\right)_{m \times n}, V_{1}\) 是 \(A X=0\) 的解空间．令 \(\alpha_{i}=\left(a_{i 1}, a_{i 2}, \cdots, a_{i n}\right)^{\prime}(i=1,2, \cdots, m), V_{2}\)为向量组 \(\alpha_{1}, \alpha_{2}, \cdots, \alpha_{m}\) 生成的向量空间，即 \(V_{2}=L\left(\alpha_{1}, \alpha_{2}, \cdots, \alpha_{m}\right)\) ，证明： \(\mathbb{R}^{n}=V_{1} \oplus V_{2}\) ．

八．设 \(M_{2}(P)\) 表示数域 \(P\) 上全体2阶矩阵关于矩阵加法和数乘构成的 \(P\) 上的线性空间，在 \(M_{2}(P)\)中定义的变换 \(\sigma\) 如下

\[
\sigma\left(\left(a_{i j}\right)\right)=\left(b_{i j}\right), \forall A=\left(a_{i j}\right) \in M_{2}(P),
\]

其中 \(b_{i j}=\left\{\begin{array}{ll}a_{i j}, & i \neq j ; \\ i \operatorname{tr}(A), & i=j .\end{array}\right.\) 这里 \(\operatorname{tr}(A)\) 表示矩阵的迹，即对角线元素之和。

\begin{itemize}

  \item[1．] 证明 \(\sigma\) 是 \(M_{2}(P)\) 上的线性变换．

  \item[2．] 求出 \(\sigma\) 的核 \(\operatorname{Ker}(\sigma)\) 的维数和一个基．

  \item[3．] 求出 \(\sigma\) 的全部特征子空间．
\end{itemize}

九。设 \(V_{1}, V_{2}\) 是 4 维欧几里得空间 \(\mathbb{R}^{4}\) 的两个子空间，其中 \(V_{1}\) 是由向量

\[
\alpha_{1}=(1,0,1,=2)^{\prime}, \alpha_{2}=(1,2,-1,0)^{\prime}, \alpha_{3}=(1,1,0,-1)^{\prime}
\]

生成的子空间，\(V_{2}\) 由向量

\[
\beta_{1}=(0,-2,2,-2)^{\prime}, \beta_{2}=(-1,3,0,4)^{\prime}, \beta_{3}=(1,5,0,2)^{\prime}, \beta_{4}=(-1,1,2,2)^{\prime}
\]

生成的子空间，求 \(V_{1}+V_{2}\) 的正交补。

十．设 \(A\) 是数域 \(P\) 上的一个 \(m \times n\) 矩阵，\(\beta\) 是数域 \(P\) 上的一个 \(m\) 维非零列向量，令

\[
W=\left\{\alpha \mid \alpha \in P^{n} \text {, 存在 } t \in P \text { 使 } A \alpha=t \beta\right\} \text {. }
\]

1．证明 \(W\) 是 \(P^{n}\) 的一个子空间。\\

2．设线性方程组 \(A x=\beta\) 的增广矩阵的秩为 \(r\) ，证明 \(W\) 的维数为 \(n-r+1\) 。
\end{document}
